% chapters/16/sub-cpmk.tex
\begin{subcpmkbox}
Bab ini menutup \CPMK{4} sekaligus menutup seluruh rangkaian materi mata
kuliah ini. Yang dibahas bukan lagi cara membuat sesuatu, melainkan cara
memeriksa apakah yang sudah dibuat benar-benar berjalan, dan cara
mempertanggungjawabkannya.

\vspace{0.5em}
\textbf{\SUBCPMK{4.3} Mengevaluasi prototipe \istilah{game} melalui
\istilah{playtesting} dan memperbaiki integrasi level serta aset} --- ranah C5,
\CPMK{4}, \CPL{04}.

\vspace{0.5em}
{\setlength{\tabcolsep}{4pt}\footnotesize%
\begin{tabularx}{\textwidth}{@{}C{1.9cm} L{4.0cm} L{2.6cm} Y@{}}
\toprule
\headrow \textbf{Indikator} & \textbf{Rumusan} & \textbf{Dibahas pada} & \textbf{Bukti} \\
\midrule
\ind{a} & Menyusun \istilah{checklist} fungsional dan lembar \istilah{playtest} &
Bagian~\ref{sec:checklist-fungsional} dan Bagian~\ref{sec:lembar-playtest} &
Checklist fungsional yang diturunkan dari GDD sendiri, dan lembar
\istilah{playtest} yang benar-benar terisi oleh penguji lain. \\
\addlinespace
\ind{b} & Mengidentifikasi \istilah{bug} dan masalah integrasi dari hasil
\istilah{playtest} &
Bagian~\ref{sec:uji-otomatis} dan Bagian~\ref{sec:log-bug} &
Catatan \istilah{bug} yang dapat diulang kembali, beserta pemeriksaan otomatis
atas level sendiri. \\
\addlinespace
\ind{c} & Menyajikan laporan perbaikan dan refleksi pemilihan teknologi &
Bagian~\ref{sec:presentasi-refleksi} &
Dokumen perbaikan yang memuat keadaan sebelum dan sesudah, serta refleksi
penggunaan Pygame dan Unity. \\
\bottomrule
\end{tabularx}}
\end{subcpmkbox}
